\documentclass[11pt]{letter}
\usepackage[a4paper,margin=2.4cm]{geometry}
\usepackage{parskip}
\usepackage[colorlinks,urlcolor=blue]{hyperref}
\pagestyle{empty}
\signature{Quang-Vinh Dang \\ Phuong-Lan Nguyen}
\begin{document}
\begin{letter}{The Editors \\ \textit{Social Networks} \\ Elsevier}

\opening{Dear Editors,}

We submit for your consideration our manuscript \textbf{``The Camouflage Ratio:
Strategic Tie Formation and the Limits of Structural Detection''} as a research
article.

The paper concerns one of the oldest uses of relational data: inferring an actor's
attributes from those of its alters. That inference rests on homophily, and the
journal has published extensively both on measuring homophily and on the conditions
under which neighbour-based inference is reliable. We ask what happens when an actor
\emph{chooses} its ties in order to defeat the inference---and we show that the
answer is not simply ``the inference degrades''.

Three results may interest your readership. First, camouflage is not a private act.
Because a tie has two ends, an actor that forms ties outside its own category
necessarily contaminates its alters' neighbourhoods in turn, and edge-endpoint
conservation fixes the rate at which it does so. The consequence is counterintuitive:
prevalence acts as protection, so the commoner such actors are, the less camouflage
each one needs. Second, we derive in closed form the point at which averaging over
alters ceases to improve on attributes alone, and show that selectively discarding
suspicious ties does not straightforwardly repair it---discarding trades bias against
the number of ties left to average over, so its benefit is bounded by the quality of
the suspicion score rather than by the discarding itself. Third, we measure the
quantity on a population of $10{,}199$ labelled Twitter accounts and find that
every recorded relation is one on which the inference is harmed, while the same
accounts behave in opposite ways across relation layers: they avoid one another where
ties are visible and converge where content is shared. The measurement is therefore a
per-relation quantity, not a property of a population---a point on which we build
directly on the journal's multiplexity literature.

We would like to be explicit about priority, since it shapes what we claim. That a
crossover of this kind exists is not our discovery: for balanced classes it is a
result of Ma et al. (2022), and the bounded interval above which relational
information becomes informative again is the second root of the same condition,
discussed by Luan et al. (2023) as the mid-homophily pitfall. Our contribution is the
asymmetric, class-imbalanced case, together with the empirical measurement. We state
this in the abstract, the introduction and the conclusion rather than leaving a
reader to discover it.

We also note two things a referee may wish to probe, and which we have tried not to
obscure. The attribute-separation parameter on which our empirical thresholds depend
is calibrated from the achievable error of a held-out linear discriminant rather than
from the moments of the attribute distributions, because the moments-based estimate is
badly inflated on high-dimensional, anisotropic attributes; we report the calibration
and its consequences. And our identity is exact in the model but requires a degree
correction on data whose classes differ in mean degree, which we document.

The manuscript is original, is not under consideration elsewhere, and has not been
published previously. All authors have approved the submission and declare no
competing interests. In accordance with the journal's double-anonymized review policy
the manuscript itself carries no author information; author details and the CRediT
statement are in the separate title page. All data analysed are third-party and
publicly available, as set out in the data availability statement; we also supply the
code needed to reproduce every reported number.

Thank you for considering our work.

\closing{Yours sincerely,}

\end{letter}
\end{document}
